\documentclass[10pt]{article}
\usepackage[a4paper,margin=1.8cm]{geometry}
\usepackage{enumitem}
\usepackage[T1]{fontenc}
\usepackage{url}
\pagestyle{empty}
\setlength{\parindent}{0pt}
\setlength{\parskip}{0.35em}

\begin{document}

\begin{center}
    {\Large\bfseries Individual Contribution Statement}\\[0.2em]
    {\large Project 1 (ODE IVP)}
\end{center}

\textbf{Name:} JIasheng Xu \hfill \textbf{Student ID:} 2251640\\
\textbf{Role this project:} Role 5 --- Testing and Validation.

\textbf{1. My specific contributions to this project:}
\begin{enumerate}[leftmargin=2em,itemsep=0.25em,topsep=0.15em]
    \item I added and migrated three edge-case validation tests in
    \texttt{Project1/test/test\_validation.py}: tighter adaptive tolerances must
    reduce global error, an oversized explicit step exposes singular-value
    crossing, and Newton failure triggers adaptive retry.  These changes are
    traceable to commit \texttt{c4ca4ee} and merge commit \texttt{b28fe39}.

    \item I independently ran both the English and Chinese entry points and the
    complete test suite, obtaining 14/14 passing tests.  I recorded the clean
    reproduction in
    \path{Project1/report_v3/checklist/VALIDATION_RUN_2026-10-03.md}
    (commit \texttt{06461a9}).  After review, I corrected the Radau comparison
    to \(1.17\times10^{-13}\) at three significant figures and regenerated and
    independently checked all 41 manifest hashes in the Windows submission
    checkout using the documented procedure (commit \texttt{838d938}).

    \item After the validation merge, I rebuilt the report PDF from
    \texttt{Project1/report\_v3/report.tex}, checked all 32 pages, and confirmed
    that Section~4 reports fourteen tests and that the contents, references,
    and appendices are complete.  I also updated the change log, validation
    checklist, and affected manifest hashes.  This report-integration work was
    merged through PR~\#3 (merge commit \texttt{e6913e6}).
\end{enumerate}

\textbf{2. The hardest conceptual obstacle I encountered:} The main challenge
was separating a local step-doubling error estimate from accumulated global
error.  I addressed it by comparing against the exact solution and turning the
expected global-error reduction into a deterministic regression test rather
than treating one accepted local step as evidence of global accuracy.

\textbf{What I would do differently:} I would begin validation from the agreed
current \texttt{main} branch, pin dependency versions, and define rounding and
manifest-hash rules before generating final evidence.  This would avoid the
path migration and last-minute rehashing required here.

\textbf{Anything I should have asked for help with earlier:} I should have asked
the report integrator to confirm the authoritative report directory and the
freeze point for final artefacts before preparing the original test PR.

\textbf{3. One piece of AI-generated output my team used, and what I changed:}
As recorded in Appendix~A, AI helped check deliverable completeness and
organise a draft of the validation evidence.  I revised that draft after manual
reproduction: I corrected the Radau rounding, distinguished close agreement
from exact matching, ran the complete suite in a clean environment, and
independently verified the documented SHA-256 results before submission.

\textbf{4. One thing I still do not fully understand:} I do not yet have a
rigorous global-error bound for this nonlinear matrix flow that starts from the
local step-doubling estimator while also accounting for rejected Newton steps.

\vfill
\textbf{Signature:} JIasheng Xu \hfill \textbf{Date:} 4 October 2026

\end{document}
